\chapter*{Agradecimentos}

Em primeiro lugar, agradeço a Nosso Senhor Jesus Cristo, Rei dos céus e da terra, e a Nossa Senhora de Fátima, que despertou em mim o interesse pela busca constante da verdade.

Tornar-me bacharel não seria possível sem o apoio incondicional de minha família, a quem, depois de Deus, devo o dom da vida. Agradeço também a Rogério, Vânia e seu filho João Pedro, a quem me une uma eterna amizade, por cederem o lar onde este sonho começou.

Aos Nerdolas, a quem devo os bons momentos da graduação: Diogo Tallys, Filipe Soriano, Hugo Coelho, Isaque Cerqueira, José Matheus, Lucas Heron, Rafael Luciano, Ricardo Vinícius, Ruan Tenório, Saulo, Thiago Ribeiro, Victor Hugo e Vinícius Neitzke. À Legião de Maria, que me deu rumo racional e espiritual nesta jornada, pelas companhias de David, João Pedro, Juliana e Klyvia, e aos colegas que me guiaram no caminho da razão: Ismael Damião e Rômulo Souza.

Ao GEDAI, primeiro grupo de pesquisa de que participei, e ao GEMA, onde fiz amizades que levo para além da graduação. Aos colegas da Edge Academy, da gestão e dos alunos, pelo ambiente de aprendizado e de trabalho em equipe, em especial a William Tcheldon, com quem trabalhei no meu primeiro projeto. À minha amiga Júlia de Oliveira, pelas boas conversas que tanto me ajudaram e apoiaram durante a graduação.

Ao meu orientador, Prof. Thales, pelo tempo e pela dedicação a este trabalho. Agradeço, ainda, aos bons docentes que se empenharam em fazer do Instituto de Computação um lugar sério e, em especial, ao professor Thiago Cordeiro e à secretária Ana Luisa, pelo empenho nos programas de extensão da universidade.

Por fim, assim como nenhum sistema se sustenta com um único componente, esta trajetória dependeu de muitas peças, várias das quais não caberiam nesta página. A cada uma delas, deixo o meu sincero agradecimento.


\begin{minipage}{.45\linewidth}
\begin{flushleft}
% Maceió, 24 de novembro de 2023.
\end{flushleft}
\end{minipage}
\begin{minipage}{.45\linewidth}
\end{minipage}